\documentclass[11pt]{article}
\usepackage[top=1in, bottom=1in, left=1in, right=1in]{geometry}

\usepackage{amsmath}
\usepackage{amssymb}

\title{Howard Math 273, Midterm Exam, \\ {\normalsize Fall 2026; Instructor: Sam Hopkins; Taken on: Wednesday, October 7th}}
\date{}

\begin{document}

\maketitle

\thispagestyle{empty}

\vspace{-1.8cm}

Each problem is worth 10 points, for a total of 40 points. You have the full 80 minute class period to work on these problems. Partial credit will be given generously, so even if you cannot fully solve a problem, show me what you know.

\begin{enumerate}

\item Define the sequence of integers $a_n$ by $a_0 = 1$, $a_1 = 2$, and $a_n = 2 a_{n-1} + 3 a_{n-2}$ for $n \geq 2$. So $a_n = 1,2,7,20,61,\ldots$ for $n=0,1,2,3,4,\ldots$. Find the real numbers $b$ and $c$ for which we have $a_n \sim c \cdot b^n$ as $n \to \infty$.

\item Define the sequence of integers $c_n$ by $c_0=1$ and $c_n = \sum_{k=0}^{n-1} c_k \cdot c_{n-1-k}$ for $n \geq 1$. So $c_n = 1,1,2,5,14,\ldots$ for $n=0,1,2,3,4,\ldots$. (These are called the \emph{Catalan numbers}.)
\begin{enumerate}
\item Form the generating function $\mathcal{C}(x) := \sum_{n \geq0} c_n x^n$. Argue, using the defining recurrence for the $c_n$, that $\mathcal{C}(x) = 1+ x\cdot \mathcal{C}(x)^2$.
\item Use the quadratic formula to conclude that $\mathcal{C}(x) = \frac{1\pm \sqrt{1-4x}}{2x}$. (We can show the $\pm$ must be a $-$ by checking the first term in the series expansion, but you don't have to do this.)
\item Show that $1-\sqrt{1-4x}=\sum_{n \geq 0} \frac{2}{n+1}\binom{2n}{n} x^{n+1}$. {\bf Hint:} You can use, without proof, that $\frac{1}{\sqrt{1-4x}} = \sum_{n \geq 0}\binom{2n}{n}x^n$ (which we showed in class), together with a derivative/integral.
\item Conclude that $c_n = \frac{1}{n+1}\binom{2n}{n}$ for all $n \geq 0$.
\end{enumerate}

\item \begin{enumerate}
\item Let $a_n$ be the number of partitions of $n$ where, for each $k \geq 1$, there are at most two parts equal to $k$. Show that $\sum_{n \geq 0} a_n x^n = \prod_{k \geq 1} (1+x^{k}+x^{2k})$.
\item Let $b_n$ be the number of partitions of $n$ where each part is congruent to $1$ or $2$ modulo~$3$. Show that that $\sum_{n \geq 0} b_n x^n = \prod_{k \geq 1} \frac{1}{1-x^{3k-2}} \cdot \prod_{k \geq 1} \frac{1}{1-x^{3k-1}}$.
\item Conclude that $a_n = b_n$ for all $n \geq 0$. {\bf Hint:} Use parts (a) and (b), and the algebraic identity $(1-x^3) = (1+x+x^2)(1-x)$.
\end{enumerate}

\item Show that, for each $n \geq 2$, the number of permutations $\sigma \in S_n$ for which $1$ and $2$ belong to the same cycle of $\sigma$ is $\frac{n!}{2}$. {\bf Hint:} Use induction. Check the base case of $n=2$. For the induction step, imagine building a permutation in~$S_n$ from one in $S_{n-1}$ by writing the permutations in cycle notation and placing the $n$ somewhere.

\end{enumerate}


\end{document}
